% subsection content: 対数型 \, $a_{n+1}=pa_n^q$

今までは$a_n$に対する操作は定数倍したり関数が掛けられたりするだけだった.
しかし,\,$a_n$に指数がついたらどうだろうか.
この場合はかなり難しい発想の変形が必要である.
それは対数を取るという操作である.
実際にやってみよう.
\[
  a_{n+1}=pa_n^q
\]
について,両辺の対数を取ると
\begin{align*}
  \log_p{a_{n+1}} & =\log_p{pa_n^q}          \\
                  & =\log_p{p}+\log_p{a_n^q} \\
                  & =q\log_p{a_n}+1
\end{align*}
$b_n=\log_p{a_n}$とおくと,
\begin{align*}
  b_{n+1}=qb_n+1
\end{align*}
あとはただの特性方程式型の問題である.
底をいくらにするかは自由だが$p$に関連する数だと扱いやすい.
たとえば$p=k^m$の時は$k$でとっておくことが多い.
作問者側は解答がシンプルな形になるようにするため,より小さい底を取ることで答えが求めやすくなる.
\noindent\textbf{【例題】}
\begin{enumerate}[label=(\arabic*),parsep=3pt]
  \item $\displaystyle a_1=2,\quad a_{n+1}=2a_n^2$
  \item $\displaystyle a_1=3,\quad a_{n+1}=9a_n^2$
  \item $\displaystyle a_1=1,\quad a_{n+1}=(n+1)a_n$
\end{enumerate}

\noindent\textbf{【方針】}
\begin{enumerate}[label=(\arabic*),parsep=3pt]
  \item 両辺の底を$2$とする対数を取り,\,$b_n=\log_2a_n$とおいて一次の漸化式に変形します.その結果,漸化式は特性方程式型になります.
  \item $9=3^2$であることに注目し,底を$3$とする対数を取りましょう.あとは(1)と同様です.
  \item 階比型に見えますが,対数型としても解けます.
\end{enumerate}

\noindent\textbf{【解答】}
\begin{enumerate}[label=(\arabic*),parsep=3pt]
  \item $\displaystyle a_n=2^{3\cdot2^{n-1}-2}$
  \item $a_n=3^{3\cdot2^{n-1}-2}$
  \item $a_n=n!$
\end{enumerate}

\noindent\textbf{【解法】}
\begin{enumerate}[label=(\arabic*),parsep=3pt]
  \item $\displaystyle a_1=2,\quad a_{n+1}=2a_n^2$\\
        両辺の底を$2$とする対数を取ると,
        \begin{align*}
          \log_2a_{n+1}
           & =\log_2{2a_n^2}           \\
           & =2\log_2{2a_n}            \\
           & =2(\log_2{2}+\log_2{a_n}) \\
           & =2\log_2{a_n}+2
        \end{align*}
        ここで
        \[
          b_n=\log_2a_n
        \]
        とおくと,
        \begin{align*}
           & b_{n+1}=2b_n+2     \\
           & b_{n+1}+2=2(b_n+2)
        \end{align*}
        ここで
        \[
          c_n=b_n+2
        \]
        とおくと,
        \[
          c_{n+1}=2c_n
        \]
        $c_1$は,
        \[
          c_1=b_1+2=\log_2{a_1}+2=\log_2{2}+2=1+2=3
        \]
        従って,
        \begin{align*}
           & c_n=3\cdot2^{n-1}       \\
           & b_n+2=3\cdot2^{n-1}     \\
           & b_n=3\cdot2^{n-1}-2     \\
           & a_n=2^{3\cdot2^{n-1}-2}
        \end{align*}
  \item $\displaystyle a_1=3,\quad a_{n+1}=9a_n^2$
        $9=3^2$\\
        底を$3$とする対数を取ると,
        \begin{align*}
          \log_3a_{n+1}
          =\log_3(9a_n^2)=2+2\log_3a_n
        \end{align*}
        ここで
        \[
          b_n=\log_3a_n
        \]
        とおくと,
        \begin{align*}
           & b_{n+1}=2b_n+2     \\
           & b_{n+1}+2=2(b_n+2)
        \end{align*}
        ここで
        \[
          c_n=b_n+2
        \]
        とおくと,
        \[
          c_{n+1}=2c_n
        \]
        $c_1=3$より,
        \begin{align*}
           & c_n=3\cdot2^{n-1}       \\
           & b_n+2=3\cdot2^{n-1}     \\
           & a_n=3^{3\cdot2^{n-1}-2}
        \end{align*}
  \item $\displaystyle a_1=1,\quad a_{n+1}=(n+1)a_n$\\
        両辺の自然対数を取って
        \[
          \log a_{n+1}=\log{n+1}+\log a_n
        \]
        したがって
        \begin{align*}
          \log a_n&=\log a_1 + \log{2}+\cdots+ \log{n}\\
                   &=\log 1 + \log{2}+\cdots +\log{n}\\
                   &=\log n!\\
                   &\therefore \quad a_n=n!
        \end{align*}
\end{enumerate}
